\todo{ensure enough complexity stuff is introduced before the next section -- enough to analyze complexities here.}

Our goal in this section is to find a method for SPV (Simple Payment Verification) between any two simplex-chains that is safe, fast, succinct, and reliable.
\begin{itemize}
    \item \emph{Safe}: an adversary must either perform a 51\% attack or break a cryptographic dependency to produce an SPV proof that appears valid but is dishonest/fraudulent (e.g., of a valid transaction that depends on a fraudulent transaction).
    \item \emph{Fast}: $O(t) \le O(\lambda) = O(1)$, where: $t$ is the maximum delay between a transaction being confirmed by a simplex-chain and the point at which a different simplex-chain begins to accept SPVs proofs of transactions in that transaction's block; $\lambda$ is a security parameter (e.g., $\lambda = 6$ local confirmations).
    \item \emph{Succinct}: for any SPV proof $\pi$, $O(\pi) < O(c)$ and $|\pi|$ is practical.
    \item \emph{Reliable}: the above properties hold for all transactions in all simplex-chains.
\end{itemize}

\todo{n/c adversary -- explain it}

Given everything we have discussed up to now, it's not clear how SPV is possible.
The problem is that miners on one simplex-chain do not validate the state transitions of other simplex-chains.
Each simplex-chain, locally, has an $O(\nicefrac{n}{c})$ share of the network's security (there are $O(c)$ many chains).
An $O(\nicefrac{n}{c})$ adversary could theoretically produce \emph{all} of the blocks for a single simplex-chain (e.g., by mining at a slight loss to push other miners out).
The adversary could then create a very long chain of blocks that have valid headers, PoRs, etc, but also contain one or more fraudulent transactions.
Perhaps the main-chain eventually corrects to a non-fraudulent history (similar to \autoref{sec:preventing-dos-attacks}) -- but in that case we'd need to wait a \emph{very} long time to be sure that such a correction \emph{must} have happened.
That violates the \emph{fast} goal, so this idea fails to meet our overall goal.

Are we stuck?
For a blockchain network to be scalable, it must be the case that miners are \emph{not} required to validate more than one chain!
Have we come all this way, attempting to deal with the core conflict of the trilemma, only to find that we just moved the conflict?

Let's not jump to that conclusion just yet -- there are cryptographic options like SNARKs and STARKs that might help.
Before that, though, let's slow down and reconsider why SPV works for Bitcoin.
Perhaps we can learn something.

\subsubsubsection{Why does SPV work for Bitcoin?}

For the sake of simplicity, let's consider Bitcoin in isolation.
Particularly, we'll assume that there are no other networks of similar size using the same PoW algorithm or compatible mining hardware.

A concise argument for the security of Bitcoin-SPV proofs is:

\begin{argumentslist}
    \item An SPV proof is a merkle branch proving that a transaction exists in a block.
    \item All blocks are part of the same blockchain and history.
    \item All transactions that occur are in some block's merkle tree.
    \item If a transaction is invalid, then the block is invalid.
    \item If a block builds on an invalid block, then it is invalid.
    \item If a miner produces an invalid block, they get no block reward.
    \item $\therefore$ Honest miners won't produce or build on invalid blocks.
    \item \label{prem:spv-7} $\therefore$ For an adversary to produce an apparently valid (but fraudulent) proof leading to a \emph{main chain} block, they'd need to out-compete honest miners (i.e., 51\% attack the network).
    \item \label{prem:spv-8} $\therefore$ Only valid transactions exist in the main chain.\footnote{
        This isn't the full picture for some blockchains (e.g., Ethereum) even though SPV still works for them -- a transaction that fails to execute can still be valid and produce a valid state transition.
        In these cases, an invalid state transition implies an invalid block, so it's customary to prove that some state has the right values rather than that a transaction exists.
    }
    \item $\therefore$ It's safe to accept SPV proofs for blocks in the main chain.
\end{argumentslist}

It appears that \cref{prem:spv-7} is our current sticking point for Simplex-SPV proofs.

But, we don't necessarily \emph{need} to replicate \cref{prem:spv-7} if we could get to \cref{prem:spv-8} by other means.

\subsubsubsection{Dialog: Error Correction I}

\aside{
    The following is a dramatized dialog based on my thoughts leading up to the afternoon of August 3\textsuperscript{rd}.
    In reality, the dialog represents several days of exploratory learning and iteration.
    Although these events have been compressed into a single dialog, it is faithful to the nature and order of significant moments and realizations.
    % That said, given the characters, you'll have to take me at my word.
}

\begin{tcolorbox}[colback=black!70,coltext=white]
    % \centering
    3\textsuperscript{rd} August 2022, mid-afternoon, UTC+10
\end{tcolorbox}

\vspace{.5em}

\dMax{Hey, time to revist SPV.}
\dSe{Great. Remind me, where are we at again?}
\dMax{
    Still stuck on existence $\implies$ validity. Proving stuff isn't the issue, it's answering: how do we know that object of the proof is part of a valid history?
}
\dMax{
    Like, each simplex-chain already knows the main chain of each other simplex-chain, but that doesn't mean much if there's one chain with an attacker making like 90\% of the blocks.
    We'd need to wait like 80 blocks for the honest branch to win out.
}
\dSe{
    Only 80 blocks?
    It's worse than that.
    If that attack can fraudulently mint new RTs and move them to another chain, then the `block rewards' for the attacker are those new RTs.
    So the attacker won't mind pumping up the difficulty to push other miners out.
}
\dMax{Oh\dots{} yeah i guess so.}
\dMax{
    Actually, i think it might even be rational for \emph{all} miners to do that.
    Not good.
}
\dMax{I'm still convinced that, if SPV is possible in other networks, then there must be a method that works for UT.}
\dSe{Hmm.}
\dMax{Any ideas?}
\dSe{
    None that you don't have already. \newline
    Why don't we try something different, see if we can brainstorm some.
}

(Note: alt section inserted here)

\dSe{How about we start by listing tools that Eth2 might use but which we don't?}
\dMax{
    Yup, okay.
    DA proofs,
    %BLS sigs,
    cryptoeconomic fraud proofs, SNARKs and STARKs, sharding / child-chains, staking, slashing, EVM, GHOST. \newline
    That's all that come to mind atm.
}
\dSe{
    Cool -- plenty for the moment. \newline
    Let's go through the list and see if we can eliminate the options one-by-one. \\
    Tell me about DA proofs.
}
\dMax{
    Certain types of Ethereum nodes will generate them to prove that the node has some data.
    The idea is that an attacker who wants to lie about having some data can't generate these proofs, and if enough of the network is sampling enough of the block data then it implies that the block is available for download.
}
\dMax{
    This doesn't help us solve the SPV problem, though.
    We already know the blocks are available because miners do simplex-wide consensus on that availability.
    Miner's must only reflect blocks that they've downloaded, and only reflect blocks that reflect blocks that they've downloaded, etc.
}
\dSe{
    Right, so we know that all the relevant blocks have been widely dispersed.
    Seems like we can rule out DA proofs for now.
    % What about BLS sigs?
}
% NOTE: this has an error: can't aggregate BLS sigs unless the sign the same thing. (thx chris)
% \dMax{
%     They're pretty neat -- mostly because of \emph{signature aggregation}.
%     If transactions had BLS sigs, then we could aggregate all of the signatures per-block into one signature.
%     It'd save space per-block, but might also make SPV stuff more complex.
%     That said, we could just use a protocol like segregated witness to achieve the same goal.
%     In any case, without introducing PoS-like consensus rules, I don't see how BLS sigs could help with SPV.
% }
\dSe{
    Okay.
    Next, Eth2's fraud proofs.
    These are different from our fraud proofs, right?
}
\dMax{
    The main difference is that we don't have any sort of staking/slashing involved.
    The general principle is the same: show that a computation was done incorrectly, or that a block contains a contradiction.
}
\dSe{How do Eth2 fraud proofs work?}
\dMax{
    Some nodes act like police -- they stake some ETH and then check that blocks/state-transitions are valid.
    If any of those police-nodes find a problem, they sign a message that challenges the invalid block.
    If the fraud proof was false, then the police-node gets slashed.
    Otherwise, whoever signed the original block gets slashed.
}
\dSe{
    Okay, so that doesn't seem compatible with PoR/UT, at least not without big changes.
    How do our fraud proofs work?
}
\dMax{
    The goal of our fraud proofs is to show only the problematic parts of the block using the minimum data necessary.
}
\dMax{
    Say an attacker made a block that looks valid in all ways except the state-transition -- e.g., there's a tx that mints 1 billion coins or w/e.
    We can always prove that such a tx is fraudulent if miners have to commit to the state-root before and after each transaction (in the block itself).
    All we have to do is provide all of the state-tree witnesses that are necessary to evaluate that one tx.
    If the tx is valid, then the output state root will always match the one recorded.
    (This sort of design is related to \href{https://github.com/zack-bitcoin/amoveo-docs/blob/master/design/stateless_full_node.md}{stateless} full nodes, which use similar structures, constraints, and patterns.)
}
\dMax{
    So, say an attacker makes an invalid block.
    The miners (and other full nodes) in that chain will try to verify the block and find the error.
    Then, instead of building on the invalid block like normal, miners would link back to that block as an invalid block along with the fraud proof.
}
\dSe{
    Right, and that only works on the local chain, so, if the attacker is making 90\% of the blocks then we're basically in the same `80 blocks' situation from before.
    Which we already know doesn't work.
}
\dMax{
    Exactly.
    We'd need to wait way too long for the situation to correct itself.
}
\dSe{
    Okay, zk stuff next.
    Btw, I've been meaning to ask:
}

% note: the next few dialog msgs overlap with some from the alt section below
\dSe{
    Why haven't you considered zk-STARKs as a useful tool for SPV?
}
\dMax{
    Well, I guess I've always associated them with layer 2 stuff, and L2 isn't really relevant to UT's consensus-level design.
    So I just put it to the side and didn't consider it.
    I ruled it out early.
}
\dSe{
    But STARKs/SNARKs are proofs of valid computation, right?
    That seems like it might be useful.
}
\dMax{
    Yeah.
    If we could produce a ``validity proof'' for each block, then that could potentially solve the SPV problem.\newline
    \newline
    Maybe each block could come with a validity proof, or one for the previous block if they're too slow to generate.
    Then, before doing an SPV proof, you could prove that the block is valid first.
}
\dMax{
    But, ofc, you'd need to prove that all the prior blocks were valid, too.
    If we did that for $O(c)$ many chains... well that's probably an issue.
}
\dSe{
    Sure, but some of those proof methods can be aggregated, right?
}
\dMax{
    True.
    Bulletproofs are one example.
    The proofs are short-ish, mb like 2-5 KB, but super slow to generate.
    Verification time is fine in practice for small proofs, but it's linear with the size of the circuit.
    So they're probably not practical without heavily rate-limiting them.
}
\dSe{
    Hmm.
}
\dSe{
    STARKs can be recursively verified, right?
}
\dMax{
    Yeah, that's sort of like aggregation.
    Basically you take some STARKs, which verify e.g., a bunch of blocks in a chain-segment, and then you produce another STARK that verifies those STARKs.
    It's actually pretty neat: the inner STARKs have complexity relative to the block-verification circuit, but the `root' STARK has complexity relative to the STARK-verification circuit.
    Since the block-verification circuit presumably includes the STARK-verification one, we'll always end up with a smaller proof.
}
\dSe{
    Sounds like there's some potential there.
    Is there a catch?
}
\dMax{
    Yeah, there is.
    The cryptographic properties of STARKs are great, but the proof size is definitely an issue.
    We're talking like 50 KB minimum, and up to 200 KB.
    There are other problems, too.
    Like that creating such proofs can take many GB of RAM and that we don't know who'd generate them.
}
\dMax{
    I guess we could have massive blocks to accommodate them, but it doesn't really feel viable.
}
\dSe{
    I agree.
}
\dSe{
    So what about SNARKs, then.
    Do they solve any of those problems?
}
\dMax{
    Yeah\dots{} sorta. \newline
    The proof size is constant, like 200-400 B depending on the scheme, and the verification time is $O(1)$ too.
}
\dMax{
    The problem is the trusted setup.
    Like, SNARKs solve lots of the problems of STARKs, but the crypto properties aren't okay.
}
\dMax{
    It's a real shame in some ways -- like the proofs are short enough and verification is fast enough that they could concievably be \emph{in the PoR graph itself}.
    Then, you'd just \emph{know} that the blocks you were doing SPV against were valid. \newline
    \dots{} actually, if we did that, we could basically enforce, \emph{at the consensus level}, that all PoRs require a SNARK.
    That'd mean that all reflected blocks are valid (since, if you were a miner, reflecting a block with an invalid SNARK would invalidate your block, and so on).
}
\dSe{
    Huh.
    That seems like a decisive advantage.
}
\dMax{
    Yeah!
    We might need to make blocks a bit bigger or adjust some parameters or something, but in principle everything would still work.
}
\dSe{
    So, talk me through the problem, then.
}
\dMax{
    Well, SNARKs have a trusted setup -- STARKs have a transparent setup so don't have this issue. \newline
    The best SNARK systems have a one-time setup where a special number gets generated.
    Then anyone else can reuse this number to prove correct execution of any suitable algorithm (though generating the proof is slow).
    The special number is the output of a special function that requires some random data as input.
    The random data is sometimes called ``toxic waste'': if it's not properly disposed of, then bad stuff happens.
    So if we used SNARKs, then the toxic waste is like an existential threat hanging over us forever -- because it's impossible to prove that it was disposed of.
}
\dSe{
    Hmm. What sort of ``bad stuff''?
}
\dMax{
    Well, most of this tech has been developed with a focus on L2 scaling.
    In that case the network operator (who also generates and publishes the proofs) could basically steal everyone's money/tokens/whatever.
    So, there's \emph{literally no way to tell} if a SNARK is honest and valid, or dishonest and just `looks' valid.
}
\dMax{
    Plus, most of these systems have some zero knowledge component as well.
    So if, e.g., account balances were not public then an attacker could just print money and no one would be able to tell.
}
\dSe{
    I can see how that'd be a big problem for confidential transactions and the like.
}
\dSe{\ldots}
\dSe{
    What if the toxic waste was well known?
}
\dMax{
    Oh, uh, I haven't really thought about it.
}
\dMax{
    I guess that some people would make valid proofs and some people would try to make fraudulent ones.
}
\dMax{
    For the bad ones, though, full nodes of that chain could always tell when the proof was bad.
    They'd find out when they validated the block.
}
\dMax{
    Oh, and they can automatically generate a fraud proof for free at that point.
    Like they already have all the witness data there to produce it.
    Then, the next time an honest miner mines a block, they'll flag the fradulent block by linking back to it as invalid and include the fraud proof.
}
\dSe{
    But we can't depend on any honest miners sticking around if the attacker's willing to mine at a loss.
    That was the original problem.
}
\dMax{
    Oh, yeah. Hmm. \newline
    If we put a SNARK in the PoR graph, why can't we put the fraud proof in it too?
    Like, \emph{any} honest full node will generate the fraud proof and auto-broadcast it.
}
\dMax{
    Then any miner of \emph{some other chain} can include the fraud proof in the PoR graph.
    (Which will protect that chain from bad SPV proofs against the bad block or its descendants.)
}
\dMax{
    We just need the fraud proof to like `pollute' the complete PoR graph.
    That's doable, because the fraud proofs are short, like $O(\log c)$ or $O(1)$ or something.
    Maybe like 500 -- 2000 bytes.
}
\dMax{
    So we need a new rule.
    Something like `if you reflect a block that contains a fraud proof for some other chain, then you need to record the fraud proof for that other chain, too, unless it's already been recorded'.
}
\dMax{
    The goal being that as soon as \emph{any} simplex-chain records a fraud proof for \emph{any other} simplex-chain, then \emph{all other} simplex-chains follow suit.
    If they don't, then they can't take advantage of those PoRs.
    But all honest miners want to do this anyway -- they want the most PoRs possible and they don't want fraud.
}
\dMax{
    And that would start happening almost immediately -- like way before the next block.
    So like within a block period or so the \emph{whole} network would know about it and have recorded it.
}
\dMax{
    An attacker would need to 51\% attack the network!!
}
\dMax{Right?}
\dSe{Hmm. Maybe.}
\dSe{Repeat the idea back, a bit more organized this time.}
\dMax{
    Yup.
}
\dMax{
    The first thing we need is for each simplex-chain to support general and short fraud proofs.
    Ideally, full nodes should automatically generate them when processing an invalid block.
}
\dMax{
    Then, each chain must support validating and recording these fraud proofs alongside PoRs.
    Sometimes, a valid block might reflect an invalid block before the miner learned of the fraud proof.
    That's okay, since we expect the fraud proof to arrive a short time \emph{after} the invalid block.
}
\dMax{
    At some point, a simplex-chain will include the fraud proof in its PoR graph (which could be any simplex-chain).
    If a miner for a different simplex-chain wants to reflect that block, then they \emph{must} include the fraud proof, also.
    If the invalid block was reflected by a past block, but the fraud proof wasn't yet known, then the miner flags the invalid block anyway, even though the corresponding PoR was in an ancestor block.
}
\dMax{
    Basically, the inclusion of these fraud proofs becomes a trigger for a fork -- blocks that don't acknowledge the fraud proof can't use PoRs via blocks that do acknowledge it.
}
\dMax{
    So, for an attacker to produce a dominant history that avoids the fraud proof, they need to do a 51\% attack.
}
\dSe{
    Hmm.
}
\dMax{
    Any criticisms?
}
\dSe{
    Not yet.
    % You'll need think about it more.
    % I'll need think about it more.
}
\dSe{
    By the way, where do the SNARKs with publicly known toxic waste fit in?
}
\dMax{
    Umm\dots{} I guess they're like a short way to verify the computations of the block, except, no that doesn't make sense b/c we need to wait and see if there's a fraud proof anyway.
    I guess they don't really do anything.
    Like, including them is basically the same as assuming a block is valid in the absence of a fraud proof, but with extra computation.
}
\dMax{So we don't need them at all.}
% note: maybe wrap up day 1 here and then have a <a few days later> banner or something.
% then put all this extra stuff -- "i've thought about it now" and on -- into the 2nd part of the dialog
\dSe{Cool.}
\dSe{
    Okay, go tell Chris.
    See what he thinks.
    }
\dMax{
    On it.
    Calling now.
    }
\dMax{\dots{} I think I just realized something about scaling.}
\dMax{
    It's not about consensus on \emph{state-transitions}. \\
    It's about consensus on the \emph{error correction} of state-transitions.
}
% \dSe{\emph{ExS\=e is silent.}}
\dSe{\dots{}}
\dSe{Does the convergence with epistemology suprise you?}
\dMax{A bit.}
\dSe{Why?}
\dMax{
    It feels like a special case of \emph{\href{https://criticalfallibilism.com/yes-or-no-philosophy-summary/}{the asymmetry between criticism and praise}}.
}
\dMax{
    IDK, I guess I didn't expect to find more convergence than we already found.
}
\dMax{
    It's exciting, but I'm unsure.
}
\dSe{
    Temper your expectations.
    There hasn't yet been enough exposure to criticism, yet.
}
\dSe{
    Don't forget, \emph{convergence} is a \emph{positive argument}.
    If it was worth anything, the inversion should be a negative argument, which would be a criticism of something.
}

\dSe{
    Give it time.
    Criticisms may yet emerge.
}
\dSe{
    Besides, we're not even done yet, you know?
}
\dMax{I know.}
\dMax{Talk soon.}
\dSe{I'll be here.}


% ########  ####    ###    ##        #######   ######       #######
% ##     ##  ##    ## ##   ##       ##     ## ##    ##     ##     ##
% ##     ##  ##   ##   ##  ##       ##     ## ##                  ##
% ##     ##  ##  ##     ## ##       ##     ## ##   ####     #######
% ##     ##  ##  ######### ##       ##     ## ##    ##     ##
% ##     ##  ##  ##     ## ##       ##     ## ##    ##     ##
% ########  #### ##     ## ########  #######   ######      #########
%
% DIALOG 2


\subsubsubsection{Dialog: Error Correction II}

\dialogHeader{11\textsuperscript{th} August 2022}

% \todo{move this to additional reasoning to a new follow-up dialog?}
\dMax{
    Hey.
}
\dSe{
    So, you've decided to explain all this through a dialog?
}
\dMax{
    Yeah.
    I was having trouble writing it out.
}
\dMax{
    A dialog felt natural once I started writing; it's helping me stay unstuck.
    I like the way that we're stepping through the process.
}
\dSe{
    Which process?
}
\dMax{
    The cycle of \href{https://criticalfallibilism.com/yes-or-no-philosophy-summary/}{\emph{conjectures and refutations}}; the process of \href{https://criticalfallibilism.com/practice-thinking-in-terms-of-error-correction/}{\emph{error correction}}.
}
\dSe{
    Why is it important to show that, rather than just the eventual solution? \\
    After all, you could treat these dialogs as another draft and replace them later.
}
\dMax{
    One of my goals is to expose these ideas to criticism, and to make that criticism as easy as possible.
}
\dMax{
    I \emph{could} just present the solution, but that obscures both my reasoning and the ideas that were rejected along the way.
    I want to know if I made any mistakes, and some of those might be procedural, so it's important to show more than just the solution.
}
\dMax{
    Writing a dialog also helps establish a track record of seeking and considering criticism;
    of thinking up criticisms and answering every one;
    of truth-seeking.
}
% \dMax{
%     How else could I convince other thinkers (especially those who may offer good criticisms) to consider my ideas?
% }
% if you can't answer the criticisms that your own mind come's up with, how will you be able to answer the criticisms that anyone else's mind come's up with?
\dSe{
    Are you prepared to be critcised on that point?
}
\dMax{
    Yes -- that's a win according to that goal.
}
\dSe{
    Very well, let's proceed.
}
\dMax{
    Have you had time to think about the SPV idea?
}
\dSe{
    Yes.
    I see two problems.
}
\dMax{
    Great.
    Let's go.
}
\dSe{
    Attackers can choose which blocks they reflect.
    So what about this case?
}
\dSe{
    An attacker first pushes all the other miners out of a chain by mining at a loss.
    Then they mine their invalid block that has some tx that will later be used for SPV fraud (maybe with intermediate transactions, too).
    The fraud proof for that block gets propagated throughout the PoR graph.
    But each other simplex-chain only records it once.
}
\dSe{
    The attacker then just skips those blocks with the fraud proof.
    Is there anything to stop the attacker using PoRs via blocks \emph{after} those blocks with the fraud proof?
}
\dMax{
    Well, a block that builds on an invalid block is itself invalid, so the other simplex-chains would just record that fraud proof, too.
    Any time the attacker generates a new block, there's a new fraud proof.
}
\dSe{
    And that fraud proof -- of a block with an invalid parent -- what is it exactly?
}
\dMax{
    It'd be like: proof2 = header2 + proof1.
    Also, some metadata to point out which ancestor, etc.
}
\dSe{
    And the fraud proof of a block with an invalid grandparent?
}
\dMax{
    I see your point -- the fraud proofs get too large (assuming they're stateless).
}
\dMax{
    Tho, hang on, if the other simplex-chains know about the first fraud proof, then they automatically know that any block that builds on the invalid block must be invalid -- no matter how long ago that was.
    No fraud proof required.
}
\dSe{
    That's assuming that the fraud proof is recorded.
}
\dMax{Yes.}
\dSe{
    So, second question: what sized bribe would prevent a miner from including the fraud proof?
}
\dSe{
    `Honest' miners can choose whether or not to record the fraud proof.
    So as long as the proof hasn't yet been recorded in any other chain's PoR graph, then there's no penalty for the miner deliberately \emph{not} including the fraud proof.
    It's basically free money for the miner.
    From the miner's point of view, someone else will include the fraud proof later on, so it only makes sense to take the money, right?
}
\dSe{
    You can't fix this by waiting longer, either, since the longer that period is, the easier it is for an honest miner to take the bribe -- there are even more future blocks that could contain the fraud proof than before.
}
\dMax{
    Well, like you said there's no pentaly for the miner, currently.
    We can, however, set a lower bound on the bribe by dividing the fraudulent block's reward between miners that record the fraud proof.
    So the attacker would need to pay at least an extra block reward per invalid block per block-period (split among all simplex-chains).
}
\dMax{
    Since the attacker loses the actual block rewards, too, the total cost after $z$ blocks is the $z$-th triangle number $+ z$.
    Approximately.
    That's $O(z^2)$ cost.
}
\dSe{
    \dots{} I'm not sure about that derivation -- something to come back to. \\
    But, now, at least, we have a basis for security parameters.
}
\dSe{
    How should the reward for including the fraud proof be split up?
}
\dMax{
    Well, we don't want honest miners to have a reason to try and game it.
    And we want all chains to record it in their PoR graph.
    My first guess is to split the reward between all chains.
    So if the block reward for the invalid block was $r$, then the miner of the first block on each chain to contain the fraud proof should get $\sim\nicefrac{r}{N_1}$.
}
\dSe{
    Hmm, we'll need to come back to those details later.
    But there are bigger problems to solve first.
}
\dSe{
    If there's something for miners to gain from producing the \emph{first} block (in that simplex-chain) to contain the fraud proof, is there an incentive for them to mine sibling blocks instead of building on the best block?
}
\dMax{
    Maybe -- only if that sibling block will appear to be the best block, though.
}
\dMax{
    The goal of such a miner is to have that sibling block be favored by the fork rule.
    That shouldn't happen with the way that we count work, though.
    Like, if the sibling block uses any PoRs via R-blocks that reflect the earlier block, then the work will end up being attributed to the earlier block rather than the common ancestor.
}
\dSe{
    Is it ever honest for a miner to act like that?
}
\dMax{Hmm. I guess not.}
\dMax{
    Actually, here's a way to tighten things up.
}
\dMax{
    Consider that miners must only reflect blocks that they've downloaded, and that condition is recursive.
    So, if a new L-block (\BL{2a}) does PoR via some R-block (\BR{1}), and that R-block reflected an L-block that isn't in the main L-block's history (\BL{2}) -- well that implies that the miner has downloaded \BL2 but has not linked back to it.
}
\dMax{
    But a miner can link back to both valid and invalid blocks.
    If \BL{2} is valid, then an honest miner should build on it, and if \BL{2} is invalid then they should flag it as invalid.
}
\dMax{
    Similarly, all blocks that \BR{1} reflected (that aren't L-blocks) should have been downloaded, so can \emph{and should} be reflected, too.
}
\dMax{
    So, we can add another two constraints to block validation:
}
\dMax{
    1. \emph{if you reflect a block, then all non-local blocks reflected by that block must also be reflected, or have been reflected by a past local block}; and
}
\dMax{
    2. \emph{you must link back to all local blocks that both are \ul[DMaxCol]{chain-tips} and have been reflected by blocks you reflect}.
}
\dMax{
    Both constraints are free, and a violation of either of these constraints should be easy to prove via a fraud proof.
}
\dMax{
    Oh, and the \emph{link back to all possible local blocks} rule would reduce the expected length of an empty-block DoS from $O({A_{\text{blocks}}}^2)$ to $O(A_{\text{blocks}})$, too.
}
\dSe{
    How does this stop miners making sibling blocks to harvest fraud proof rewards?
}
\dMax{
    Say that the local block \BL{2} contains the fraud proof for some other chain.
    A miner could make a block, \BL{2a}, that looks like it was created earlier than \BL{2}.
    They'd just set the timestamp earlier and only include some of the PoRs that \BL{2} has.
    However, there are PoRs from other chains that point to \BL{2}, but not \BL{2a}.
    So the next L-block (\BL{3}) will link back to both, but \BL{2} will be the primary parent.
}
\dMax{
    If a miner is mining \BL{2a} after \BL{2}, they're now unable to include any PoRs via any blocks that reflect \BL{2} (and this applies recursively).
    So there's no way for \BL{2a} to have substantially more PoRs than \BL{2}.
    There might be \emph{some} PoRs that become available after \BL{2} was mined but before the miner downloaded \BL{2}.
}
\dMax{
    But from \BL{3}'s perspective, those PoRs are valid for both \BL{2} and \BL{2a}.
    And there will be some PoRs for \BL{2} that don't include \BL{2a} -- the PoRs produced after other miners learn of \BL{2} but before \BL{2a} was mined.
}
\dMax{
    So \BL{2} should always out-weigh \BL{2a} $\implies$ \BL{2} always wins.
}
\dSe{
    I see.
}
\todo{Check that we cover safe and reliable explicitly (or add)}
\dSe{
    Let's take stock for a moment.
    We've covered \emph{safe} and \emph{reliable}.
    What about \emph{succinct}?
}
\dMax{
    The SPV proof is a merkle/verkle branch leading to some transaction or state, which is $O(\log c) + O(1)$.
    Blocks will have room for $\sim 100$ transactions, so verkle branches could potentially be $O(1)$.
    Transactions involving SPV proofs will be larger than average, but small enough that many could be included in each block.
}
\dSe{
    Okay.\\
    What about \emph{fast}?
}
\dMax{
    I'm not sure yet, but I have an idea for how to figure it out.
}
\dSe{
    Via analyzing bribes?
}
\dMax{
    Yup.
    The goal is finding a way to tell if it's safe to accept an SPV proof.
}
\dMax{
    We know that an attacker would need to bribe all the miners for some number of blocks to pull off some fraud, and we can calculate that cost.
    The attacker's goal is to profit, so for some fraudulent transaction with value $v$, and \BribeCost spent on bribes, the attacker's profit will be $\approx v - \BribeCost$.
    And \BribeCost depends on how long the attack lasted.
    So if the attack lasts $z$ blocks, and \BribeCost is a function of $z$, then it should be safe to accept the SPV proof when $v < \BribeCost(z)$.
}
\dSe{
    When does $z$ start?
}
\dMax{
    When the fraudulent transaction is made.
}
\dSe{
    Can't the attacker parallelize that to reduce the waiting period?
    They only have to pay $\BribeCost(z)$ \emph{per-chain}, not \emph{per-transaction}.
}
\dMax{
    Ahh, yeah.
}
\dSe{
    And how do you know which txs will later be used for SPV?
}
\dMax{
    We don't, yet -- b/c the target chain is the \emph{outgoing} chain. \\
    I think we can solve both problems at once, tho.
    First, we can flag all txs that are meant as outgoing SPV transactions.
    That tells us which txs matter.
    Second, we can measure the waiting time for all the relevant txs and run them in batches.
}
\dSe{
    How does this batch system work?
}
\dMax{
    Each chain could track outgoing transactions into a batch and commit to that via a field in the header.
    Each batch has its own ID number, too.
    As txs accumulate, the total outgoing value increases until it hits the $v < \BribeCost(z)$ breakpoint.
    Then the counter resets to 0.
    Other simplex-chains require that an SPV proof links back to a transaction that's part of a batch that is canonical according to the main chain and that the batch has a large enough $z$.
}
\dMax{
    If $O(\BribeCost(z)) = O(z^2)$ -- like I guessed before -- then $O(z) = O(\sqrt{\BribeCost}) = O(\sqrt{v})$.
    So the waiting time is roughly proportional to the square root of the value being moved.
}
\dSe{
    It would be troubling if the waiting time didn't scale like that.
}
\dSe{
    But if that's the case, then can't the attacker exploit that complexity by adding a late transaction, close to the batch execution?
}
\dMax{
    Yes\dots{} that bad tx might only have a few confirmations, so $v < \BribeCost(z)$ isn't upheld.
}
\dMax{
    We could fix that by adding transactions to the \emph{next} batch, not the current one.
}
\dMax{
    Then we can guarentee that \emph{all} outgoing transactions have been through some minimum waiting period, $z$.
    Somewhere between $z$ and $2z$ blocks if batches are regular.
}
\dSe{
    Hmm, okay.
}
\dSe{
    Let's return to \BribeCost.
    Can you derive it?
}
\dMax{
    Maybe, the idea isn't well formed yet, though.
}
\dSe{
    Then let's discuss it a bit.
    See if we can come up with some understanding of how it behaves.
}
\dSe{
    What's your intuition so far?
}
\dMax{}

\begin{comment}
    - bribe all miners
    - 1 tx per chain
    - miners unpredictable
    - needs to be most profitable option
    - bribe > opportunity cost
    - op cost is sum of rewards of that chain's share of invalidated blocks
    - op cost is proportional to z
    - bribe needs to be paid z times
    - so ~z^2
\end{comment}

-------------------

\dSe{
    You know, if the waiting time for SPV transactions is $O(\sqrt{v})$, that doesn't the \emph{fast} criterion.
}
\dMax{
    I know.
}
\dSe{
    Have you considered what to do in that case?
}
\dMax{
    Yes.
    Discuss it.
}
\dSe{
    And the \emph{fast} criterion?
}
\dMax{
    Leave it be.
    If we need to change our goals, then we'll address that when the time comes.
}

% \dMax{
%     I think we need to go back to the cost of bribes.
% }
% \dMax{
%     Consider the targeted simplex-chain (the one that would have an `outgoing' tx).
%     Also assume the attacker is the only miner on that chain.
% }
% \dMax{
%     For a bribe over $z$ blocks, the attacker's cost ($Z$) is $Z = r(\frac{1}{2} z (z+1) + z) = \frac{1}{2} r z (z + 3)$.
% }
% \dMax{
%     If the total value of outgoing SPV transactions (for that block) have value $v$, then an attack is profitable if $2v > rz(z+3)$.
%     We want the opposite inequality for some idea of when an SPV transaction is safe to accept.
% }
% \dMax{
%     So, $2vr^{-1} < z(z+3) = zz + 3z$.
% }

\todo{discussion and investigation of bribe attacks and various options / configs}

\todo{branch: bribes when no FP recorded, bribes when an FP recorded.}


todo:

\begin{itemize}
    \item why do we only care about outgoing SPV?
    \item bribe threshold discussion
    \item - only pay 1/2 miners?
    \item - can briber profit?
    \item - 'treasury' for fraud proofs? -- attacker can probs profit
    \item - what about 1/sqrt(n)
    \item - security parameter, z local confs and N*z global work?
    \item nipopows
    \item accelerate via staking?
    \item fraud spv transactions need to be made during the batch prior to that tx's batch. so the fraud tx occurs \emph{before} the first block that txs in the \emph{current} batch can be used.
    \item add rule: spv to an invalid block invalidates your block.
    \item thus: spv txs to the current batch will invalidate remote chains. (not a problem b/c all nodes know there's a fraud proof, even if miners are bribed. we just need to ensure users can require that a particular FP was recorded.)
    \item so miners don't want to do SPV against those blocks b/c it'll invalidate their bribe.
    \item the chain needs to remain valid for the bribe to work.
\end{itemize}


% \begin{align}
%     0 &= zz + 3x - 2v/r \\
%     z &= -3/2 \pm \sqrt{9 + 8v/r}/2 \\
%     z &= (-3 + \sqrt{9+8v/r})/2 \\
%     \text{if } v < r \\
%     z &< (-3 + \sqrt{17} < \sqrt{25})/2 \\
%     z &< (-3 + 5)/2 \\
%     z &< 1 \\
%     \text{if } v < xr \\
%     2z &< -3 + \sqrt{9 + 32 = 41 < 49} \\
%     2z &< -3 + 7 \\
%     z &< 2
%     \text{if } v < xr \\
%     2z &< -3 + \sqrt{9 + 8x = 1 + 8(1+x)} \\
%     2z &< -3 + 2\sqrt{2(9/8 + x)} \\
%     2z &< -3 + 2\sqrt{2}\sqrt{2x} \\
%     2z &< -3 + 4\sqrt{x} \\
%     z &< 2\sqrt{x} - \frac{3}{2}
% \end{align}

spv via bulletproofs?
or something similar?
basically the job is to prove that all inputs, since the last bulletproof, are valid.
like you JUST prove that -- not all the blocks, just the chain of inputs

if > some size -- makes sense mb?

% ########  ########  #### ########  ########
% ##     ## ##     ##  ##  ##     ## ##
% ##     ## ##     ##  ##  ##     ## ##
% ########  ########   ##  ########  ######
% ##     ## ##   ##    ##  ##     ## ##
% ##     ## ##    ##   ##  ##     ## ##
% ########  ##     ## #### ########  ########
%
% Bribe
%
% 88""Yb 88""Yb 888888    db    88  dP 88""Yb  dP"Yb  88 88b 88 888888 .dP"Y8
% 88__dP 88__dP 88__     dPYb   88odP  88__dP dP   Yb 88 88Yb88   88   `Ybo."
% 88""Yb 88"Yb  88""    dP__Yb  88"Yb  88"""  Yb   dP 88 88 Y88   88   o.`Y8b
% 88oodP 88  Yb 888888 dP""""Yb 88  Yb 88      YbodP  88 88  Y8   88   8bodP'
%
% Breakpoints
%

\subsubsubsection{SPV Bribe Attacks: Derivations \& Breakpoints}

\todo{derive bribe attack threshholds}



Let's assume that an attacker is willing to bribe miners to exclude a fraud proof.
The goal of the attacker is to continue this attack beyond whatever , how much will that cost them over the

First, let's establish a baseline: what does a healthy simplex look like?
\emph{Healthy} meaning: no fraud proof is possible because all blocks and state-transitions are valid.

We need to know what a healthy simplex (and its PoR graph) looks like because a successful attack will be indistinguishable.
Individual nodes of the network may all be aware of a fraud proof when an invalid block is mined -- and we expect this if the P2P layer is working correctly.
However, during an ongoing (successful) attack, that fraud proof is not recorded in the PoR graph.
The fraud is `invisible' when considering just the blockchain data itself, which means that SPV proofs may be fraudulent.

qualities:
- maximal interconnection (ish) in PoR graph;
  all miners reflect all other chains and link back to all other blocks of their own chain.

- some 'fuzziness', though, where blocks were created near-simultaneously.
  these are the only cases where blocks shouldn't reflect one another or build on other L-blocks.
  (excluding cases like a bug which causes a network fork; e.g., bitcoin march 2013.)

When fraud is detected, the PoR graph goes through multiple phase changes until normality is restored.

\begin{center}
    \tikzstyle{n} = [align=center, draw, rectangle, rounded corners, inner sep=5pt]
    \tikzstyle{ar} = [->, anchor=north, rounded corners, shorten >= 2pt, shorten <= 2pt]
    \tikzstyle{label} = [align=center, anchor=center]
    \begin{tikzpicture}[node distance=1.8cm]
        % states of PoR graph
        \node (s0) [n] {0. PoR Graph \\ Looks Healthy};
        \node (s1) [n, right=of s0, xshift=.95cm] {1. Attack \\ In Progress};
        \node (s2) [n, right=of s1, xshift=.6cm] {2. Partial FP \\ Acknowledgement};
        \node (s3) [n, below=of s1.north, yshift=.2cm] {3. Complete FP \\ Acknowledgement};

        % group s0 and s1
        \node[n, fit=(s0) (s1), dashed, inner sep=7pt] {};

        % arrows and labels
        \draw [ar] (s0) -- node [label] {Fraudulent \\ block mined} (s1);
        \draw [ar] (s1) -- node [label] {FP first \\ recorded} (s2);
        \draw [ar] (s2) -- (s2 |- s3) -- node [label] {FP `pollutes' \\ PoR graph} (s3);
        \draw [ar] (s3) -- node[label] {PoR graph now \\ resynchronized} (s3 -| s0) -- (s0);
    \end{tikzpicture}
\end{center}

The dashed outline around phase 0 and 1 indicates that there is no sign of the fraud within the PoR graph, yet.
Nodes connected to the P2P network should be aware of the fraud, but no chain has yet published a block with the FP.

The attacker's goal is to keep the network from reaching phase 3.
Ideally, the network never transitions from phase 1 to phase 2.

The attacker and miners (both bribed and non-bribed) experience different dynamics during each of the phases.

For example, there is a qualitative difference between phase 1 and phase 2 and it is in the attacker's interest to prevent the network transitioning to phase 2 if possible.
Therefore it is most valuable for the attacker if they can prevent the FP being recorded.

Once the network has reached phase 2, it is only a matter of time before the network reaches phase 3, at which point we are back at the start.

If the network is kept in phase 1, then the PoR graph and chain data is indistinguishable from a healthy PoR graph.



\subsubsubsubsection{Phase 1 $\not\to$ Phase 2 on the Target Chain}

Reward for honest block is an extra $z/N_1$ of the block reward.
So after $N_1$ blocks the block reward has doubled.

A new honest block thus becomes heavily incentivized over time.

To mitigate this, the attacker must raise the difficulty such that all other miners find it comparatively worse than mining other chains.

Thus the attacker suffers losses compared to mining honestly: $Hkr$ of the block reward for the k-th block (e.g., $H = \frac{1}{N_1}$).
That reward is lost to work securing the target chain that would otherwise earn the miner that much revenue.
Additionally, the miner loses the block rewards each time -- so this is $r(1 + Hk)$ at round $k$.
Let $r \cdot l_m$ represent the total losses due to mining: factoring out $l_m$ in this way allows us to analyze the situation independent of the block reward.

\begin{align}
    l_m &> \sum\limits_{k=1}^z \Bigl( 1 + Hk \Bigr) \nonumber \\
    &= z + H \sum\limits_{k=1}^z k \nonumber \\
    &= z + H \frac{z (z + 1)}{2}
    \label{eq:spv-atk-losses-mining}
\end{align}



\subsubsubsubsection{Phase 1 $\not\to$ Phase 2}

\begin{itemize}
    \item Each miner (on each simplex-chain) needs bribing (but only if they make a block).
    \item Each round, the attacker must pay $rZ_k$ in bribes (where $k$ is the round, starting at 1).
    \item These bribes are split between $N_1 - 1$ miners.
    \item A miner receives $Hkr$ coins as a reward for publishing a fraud proof that invalidates $k$ blocks.
    \item Each round, miners can choose to publish the fraud proof, which earns them $Hkr$ in rewards.
    \item Thus, each round, miners have an opportunity cost of $Hkr$ if they do \emph{not} take the bribe.
    \item This only makes rational sense if the miner's bribe is greater than their opportunity cost: $rZ_k / (N_1 - 1) > Hkr$; \emph{and} the miner is \emph{guaranteed} to receive the bribe, regardless of whether the attack succeeds.
    \item The attack lasts $z$ rounds.
    \item Let $r \cdot l_b$ be the total losses due to bribes.
\end{itemize}

Therefore, we have:
\begin{align}
    % \text{let } H(z) &= \frac{z}{N_1} \nonumber \\
    % \implies N_1 H(z) &= z \nonumber \\
    Z_1 &> H (N_1-1) \nonumber \\
    Z_2 &> 2 H (N_1-1) \nonumber \\
    Z_k &> k H (N_1-1) \nonumber \\
    \text{let } l_b &= Z_1 + \cdots + Z_z \nonumber \\
    \therefore l_b &> H (N_1-1) \sum\limits_{k=1}^z k \nonumber \\
    l_b &> H (N_1-1) \frac{z (z+1)}{2}
        \label{eq:spv-atk-losses-bribes}
\end{align}

Combining \autoref{eq:spv-atk-losses-mining} with \autoref{eq:spv-atk-losses-bribes} yields:
\begin{equation}
    l_z = l_b + l_m > z + H N_1 \frac{z(z+1)}{2}
    \label{eq:spv-atk-losses-total}
\end{equation}

\subsubsubsubsection{Phase 1 $\to$ Phase 2}

The network transitions from phase 1 to phase 2 as soon as a first miner produces a block that includes the fraud proof.

- attacker
- dishonest nodes: the ones that take the bribe and help the attacker
- honest nodes: the ones that produce blocks where the FP is included in the PoR graph

Once this happens, an asymmetry emerges: honest miners (which include the fraud proof) produce proofs of reflection that are \emph{not} usable by the dishonest miners.
If the dishonest miners tried to use those PoRs, they'd produce invalid blocks unless they also include the FP.
Thus, using those PoRs effectively turns dishonest miners into honest miners.



\subsubsubsubsection{Phase 2 $\to$ Phase 3}


The duration of this transition depends on the proportion of honest miners ($h$, where $h + d = 1$ and $d$ is the proportion of dishonest miners).

Similar to \autoref{sec:dos-and-dags}, the asymmetry between honest and dishonest miners means that, even for a small $h$, the chain-weight of the honest chain-tips will dominate the chain-weight of dishonest chain-tips.
At that point, each simplex-chain's main chain will include the fraud proof, and the dishonest chain-tips will never be able to out-weigh the honest ones.

Consider the case where only 1\% of miners refuse the bribe ($h = 0.01$), and that the other 99\% of miners continue to mine only blocks that do not include the fraud proof.

Let us represent the weight of the honest and dishonest chain segments with a vector.
This is purely for convenience so that we can analyze both at the same time.
The top row is the weight of the honest chain segment, and the second row is the weight of the dishonest chain segment.
At the start of the attack ($t_0$), neither chain segment has any blocks, so $t_0 = \vec{0}$.

Each round, we expect that $h$ proportion of new blocks were honest, and $d$ proportion are dishonest.
So, after 1 block period, we expect:
\begin{equation*}
    t_1 = t_0 + \vecXY{N_1 h}{N_1 d} = N_1 \vecXY{h}{d}
\end{equation*}

After 2 block periods, the $h$ proportion of new blocks that are honest can use PoRs via the dishonest blocks.
Due to the asymmetry, the dishonest blocks can not.
\begin{equation*}
    t_2 = t_1 + \vecXY{N_1 (h+d)}{N_1 d} = N_1 \vecXY{h}{d} + N_1 \vecXY{1}{d} = N_1 \vecXY{(h + 1)}{2 d}
\end{equation*}

Each round, the honest chain segments grow by the weight of $N_1$ blocks, but the dishonest chain segments grow by $N_1 d$ blocks.
Therefore, at the $k$-th round, $k \ge 2$:
\begin{align}
    t_k &= t_{k-1} + N_1 \vecXY{1}{d} = t_{k-2} + 2 N_1 \vecXY{1}{d} \nonumber \\
    \text{observe: } t_k &= t_{k-i} + i N_1 \vecXY{1}{d} & \text{for some } i \in [1, k) \nonumber \\
    \intertext{set $i = k - 1$}
    t_k &= t_1 + N_1 (k-1) \vecXY{1}{d} \nonumber \\
    t_k &= N_1 \vecXY{h}{d} + N_1 (k-1) \vecXY{1}{d} \nonumber \\
    \therefore t_k &= N_1 \vecXY{h + k - 1}{dk} \label{eq:h-to-d-chain-segs}
\end{align}

Thus, the honest chain-tips dominate when
\begin{align*}
    h + k - 1 &> dk \\
    k - dk &> 1 - h \\
    k(1 - d) &> 1-h \\
    h k &> 1-h \\
    k &> h^{-1} - 1
\end{align*}

If $h = \frac{1}{N_1}$, i.e., for each honest miner there are $N_1 - 1$ dishonest miners, then this inequality becomes:
\begin{equation*}
    k > N_1 - 1
\end{equation*}



% ########  ####    ###    ##        #######   ######       #######
% ##     ##  ##    ## ##   ##       ##     ## ##    ##     ##     ##
% ##     ##  ##   ##   ##  ##       ##     ## ##                  ##
% ##     ##  ##  ##     ## ##       ##     ## ##   ####     #######
% ##     ##  ##  ######### ##       ##     ## ##    ##            ##
% ##     ##  ##  ##     ## ##       ##     ## ##    ##     ##     ##
% ########  #### ##     ## ########  #######   ######       #######
%
% dialog 3

\subsubsubsection{Dialog: Error Correction III}

--------

ex se: if you can't answer the criticisms that your own mind come's up with, how will you be able to answer the criticisms that anyone else's mind come's up with?

--------

in response to writing a dialog and honestly and fully exploring everything.

\dSe{(draft) if you can't answer the criticisms that your own mind come's up with, how will you be able to answer the criticisms that anyone else's mind come's up with?}


% end of day 1
% to go at the end:

\dMax{I think I just realized something about scaling.}
\dMax{
    It's not about consensus on \emph{state-transitions}. \\
    It's about consensus on the \emph{error correction} of state-transitions.
}

% note: start 2nd part/day here?


--------



\subsubsubsubsection{alt section in dialog}

\todo{
    I removed this section from the main dialog b/c it is maybe innapropriate, and a little exaggerated.
    But it ties into the theme of the dialog pretty well.
}

\dSe{Tell me about Ethereum's scaling tech; maybe they've come up with some ideas that are useful.}
\dMax{I don't think Eth2's sharding or SPV methods are going to be very useful for us, it's all just cryptoeconomic game theory.}
\dSe{
    \dots{} Max, slow down.
    Way too many errors just there to make progress.
}
\dSe{
    If we don't fix them now, we'll be \href{https://criticalfallibilism.com/overreach-summary/}{overreaching}.
    So, time to do some error correction.
    Tell me where you went wrong.
}
\dMax{Yeah, you're right. Okay.}
\dMax{
    I can see 4 errors.
    I didn't respond to what you asked;
    I had a negative bias and assumed those ideas wouldn't help;
    it's not just cryptoeconomics, there are other things like BLS signatures and zk-rollups;
    and the cryptoeconomics comment was a bit rude/dismissive.
}
\dSe{
    Right.
    All those things are related.
    We can avoid those mistakes by slowing down; by paying more attention to our biases and emotions.
    Don't let them cloud your thinking.
    If you get defensive, it'll shut down error correction.
    Consciously maintain a truth-seeking mindset.
    Remember: criticism is a gift.
}
\dSe{
    There's a big one that you missed. \newline
    Consider our goal right now, and what might stop us realizing it?
}
\dMax{
    Well, our goal is brainstoming -- ahh, yeah, \href{https://criticalfallibilism.com/brainstorming-advice/}{suppression}.
}
\dSe{
    Right. \newline
    So, tell me, why haven't you considered zk-STARKs as a useful tool for SPV?
}
\dMax{
    Well, I guess I've always associated them with layer 2 stuff, and L2 isn't really relevant to UT's consensus-level design.
    So I just put it to the side and didn't consider it.
    I ruled it out early.
}
\dSe{
    See how suppression might be holding us back? \newline
    Okay, let's get back to brainstorming.
    How about we start by listing tools that Eth2 might use but which we don't?
}


\subsubsubsection{SPV Limits}












\subsubsubsection{What about SNARKs and STARKs?}

SNARKs and STARKs, and their zk- siblings, refer to cryptosystems (or their respective proofs) with some remarkable properties.
Particularly, these systems produce proofs of \emph{correct computation} for a known algorithm, and \emph{the proof can be verified both faster and using less data than doing the computation itself}.
Even more remarkable is that large or private inputs to the computation can be hidden, and hiding inputs has negligible impact on the size of the proof and the verification time.

Vitalik wrote\footnote{
    Note that the quote is in the context of zk-SNARKs particularly, but the same idea applies to zk-STARKs, too.
}:

\bquote{
    For example, you can make a proof for the statement "I know a secret number such that if you take the word `cow', add the number to the end, and SHA256 hash it 100 million times, the output starts with \texttt{0x57d00485aa}".
    The verifier can verify the proof far more quickly than it would take for them to run 100 million hashes themselves, and the proof would also not reveal what the secret number is.
}[\href{https://vitalik.ca/general/2021/01/26/snarks.html}{An approximate introduction to how zk-SNARKs are possible / vitalik.ca / 2021-01-26}]

This is good motivation to explore whether any of these cryptosystems are useful for constructing a suitable Simplex-SPV method.

\subsubsubsubsection{zk-STARKs}

The major reason to consider S\textbf{T}ARKs first is their \textbf{t}ransparent (trustless) setup.

At first glance, zk-STARKs seem promising: complexity of the proof size is $O(\log^2 |C|)$ for an algorithm with size $|C|$, and verification time is $O(\log |C|)$.\footnote{
    These values and complexities are taken from \href{https://www.youtube.com/watch?v=h-94UhJLeck&t=1519s}{ZK Whiteboard Sessions - Module One: What is a SNARK? by Dan Boneh}.
}
(The algorithm, $C$, would be the block-verification algorithm including state-transitions.)

However, the \emph{actual} proof size is more like 50 KB in the best case.
This is problematic unless simplex-chain block sizes are much greater than 50 KB.
Realistically, we should expect a single simplex-chain to record and verify many of these proofs -- as many as one for each block in each other simplex-chain.
There are techniques to help here -- like \emph{recursive} STARKs -- but making the method more complex exacerbates other problems.
For example: creating a STARK proof is somewhat resource intensive.
The running-time is okay, but RAM requirements can be in the many-gigabytes.
 these proofs would need to be verified for each block in each other simplex-chain that

\begin{itemize}
    \item transparency, good
    \item proof size, bad
    \item maybe combining proofs?
    \item realistically, would be hard to make this feasible
    \item come back to it if we run out of other options
\end{itemize}

\subsubsubsubsection{zk-SNARKs}

\begin{itemize}
    \item trusted set-up / toxic waste
    \item why toxic waste bad
    \item what if the toxic waste was public
    \item what's the problem (fraudulent proofs)
    \item that is a problem with ZK
    \item but we can provide fraud proofs\dots
    \item wait, if we can provide fraud proofs (and we'd rely on that), why do we need a zk-SNARK?
\end{itemize}

\subsubsubsection{What about Fraud Proofs?}




\subsubsubsection{draft content below}


It might be feasible for someone to create an invalid block that gets reflections.
We expect this because PoR does not validate the state transitions of reflecting chains.
So, just because some transaction exists, or some state can be proven, doesn't mean that we can use that proof.

There's another case where we can't rely on SPV proofs: blocks that are not part of a block-DAG's main chain (since those transactions are eventually executed in a different context than the block was created in).

We know -- from \autoref{sec:dos-and-dags} -- that an attacker with more than 50\% of a chain's local hash-rate can produce relatively long chain-segments of empty (or invalid) blocks.
We also know that there will be some delay before honest miners' blocks are included in the main-chain once again.

If an attacker on some simplex-chain is producing empty-but-valid blocks, that doesn't compromise SPV.

----------------

If you have followed everything up to this point and have been considering how $\UT1$ might be implemented, you have probably noticed an omission that we have not discussed: SPV.

Let's consider:
\begin{enumerate}
    \item How does SPV work in Bitcoin?
    \item Why wouldn't SPV work in UT?
    \item How could we facilitate SPV in UT?
    \item What is actually required to do SPV in UT?
\end{enumerate}

--------

exists => valid -- technicality when txs can fail, but we just check the state for the success of that transaction instead.

--------

don't do consensus on state transitions; do consensus on error correction of state transitions.

epistemic asymmetry between praise and criticism.

---------

limit of proofs of invalididty are always smaller/easier than contraproofs of fraud
